
\begin{figure}[t]
    \centering
    \includegraphics[scale=0.46]{figure/intro.pdf}
    \caption{An illustrative example of using Qwen2.5-7B to generate a SQL test case for DBMS.}
    \label{fig:intro_case}
\end{figure}

    
\noindent Database Management Systems (DBMSs) are fundamental infrastructure for modern data-driven society, underpinning applications across finance, healthcare and e-commerce~\cite{stonebraker1986design,raasveldt2019duckdb,ozsu1996distributed}. Ensuring the correctness and robustness of DBMSs is vital, as even minor faults can lead to severe consequences, including data corruption, service outages, or security breaches~\cite{lo2010framework,chopra2010database,pick2025checking}. An essential step toward assuring DBMS reliability is thorough testing using high-quality SQL test cases, which help uncover subtle defects before deployment and prevent them from affecting production environments~\cite{zhong2020squirrel,bruno2006generating}. Traditional approaches for generating test cases, such as BuzzBee~\cite{yang2024towards}, SQLsmith~\cite{SQLsmith}, SQLancer~\cite{rigger2020testing}, can be effective for a specific DBMS; however, there exist many DBMSs designed for different application scenarios, such as DuckDB~\cite{raasveldt2019duckdb} for analytical workloads, PostgreSQL~\cite{stonebraker1986design} for general-purpose OLTP systems, and SQLite~\cite{owens2010sqlite} for lightweight embedded environments. Adapting these methods to different DBMSs requires substantial manual effort, such as crafting grammar rules and maintaining specific operators, which severely limits their scalability~\cite{zhong2025scaling,zhong2025testing,zhong2020squirrel}.
Recent advances in Large Language Models (LLMs) present promising opportunities for DBMS test case generation. First, LLMs have already demonstrated success in related data management tasks, including text-to-SQL~\cite{fan2024combining,li2024dawn,li2023can}, data cleaning~\cite{li2024autodcworkflow,naeem2024retclean,fan2024cost,qian2024unidm}, knob tuning~\cite{lao2023gptuner,giannakouris2025lambda}, and DBMS diagnosis~\cite{zhou2023d,zheng2024revolutionizing,singh2024panda}. Second, LLMs are applied in testing across multiple domains, such as general software testing~\cite{wang2024software,xue2024llm4fin}, compiler testing~\cite{gu2023llm,xie2025kitten}, and embedded testing~\cite{bley2025protocol}. As illustrated in Figure~\ref{fig:intro_case}, LLMs can directly generate SQL test cases by leveraging their understanding of SQL syntax and semantics acquired during pre-training. Despite these encouraging developments, deploying LLMs for industrial-scale DBMS testing still faces two critical challenges:

\begin{itemize}[leftmargin=*]

\item  \textbf{Challenge 1: Limited model adaptability in proprietary and resource-constrained industrial DBMS environments.} 
Industrial DBMSs often have proprietary SQL dialects with unique syntax and features. At the same time, security and privacy requirements restrict organizations to locally deployed lightweight LLMs (typically 70B parameters or smaller). The limited model capacity results in generating overly simple queries with insufficient testing coverage (see Figure~\ref{fig:challenges}(a)). Moreover, these models lack sufficient training on proprietary SQL variants, often producing queries that violate dialect-specific rules (see Figure~\ref{fig:challenges}(b)). As a result, adapting lightweight LLMs to proprietary industrial DBMSs remains challenging.

\item \textbf{Challenge 2: Limited exploration of deeper execution paths in test generation.}
Effective DBMS testing requires queries that cover meaningful execution paths beyond syntactic correctness. However, simply generating a large number of test cases does not translate into high coverage, as many queries are semantically similar and exercise only shallow execution paths. As demonstrated in a prior study~\cite{wang2021industry}, the generated test cases may improve coverage initially, but the coverage quickly plateaus as new inputs rarely explore deeper logic. Establishing an effective mechanism to guide test generation toward unexplored code regions remains a critical challenge.

\end{itemize}

\begin{figure}[t]
% \centering
% ========== Subfigure (a) ==========
\begin{subfigure}[t]{0.38\textwidth}
\begin{minipage}{\linewidth}   % 这里一定用 \linewidth！！
\begin{lstlisting}[language=SQL, basicstyle=\ttfamily, frame=single,  breaklines=true]
CREATE TABLE products (
    id INTEGER,
    name TEXT,
    price REAL
);
INSERT INTO products VALUES 
            (1, 'Laptop', 999.99),
            (2, 'Phone', 499.99);
SELECT name, price 
FROM products 
WHERE price > 10;
\end{lstlisting}
\end{minipage}
\caption{Simple query with limited testing coverage.}
\label{fig:challenge_semantic}
\end{subfigure}
% }
% \hfill
% ========== Subfigure (b) ==========
\begin{subfigure}[t]{0.38\textwidth}
% \centering
\begin{minipage}{\linewidth}
\begin{lstlisting}[language=SQL, basicstyle=\ttfamily, frame=single, escapechar=@]
CREATE TABLE employees (
  id INTEGER,
  details @\textcolor{red}{STRUCT<age: INTEGER>}@
);
INSERT INTO employees VALUES (1, 5);

SELECT * 
FROM products 
WHERE name @\textcolor{red}{REGEXP\_MATCHES}@('^A')
@\textcolor{red}{CROSS JOIN}@ cities c @\textcolor{red}{ON}@ c.id=1;
\end{lstlisting}
\end{minipage}
\caption{Query violating syntax rules on DuckDB (errors highlighted in \textcolor{red}{red}).}
\label{fig:challenge_syntax}
\end{subfigure}
% \hfill
% ========== Subfigure (c) ==========
% \begin{subfigure}[t]{0.4\textwidth}
% \centering
% \begin{minipage}{\textwidth}
% \begin{tikzpicture}[scale=0.8]
% \begin{axis}[
%     xlabel={Number of Test Cases},
%     ylabel={Branch Coverage},
%     xmin=0, xmax=1500,
%     ymin=0, ymax=50000,
%     xtick={0,300,600,900,1200,1500},
%     ytick={0,10000,20000,30000,40000,50000},
%     legend pos=south east,
%     ymajorgrids=true,
%     grid style=dashed,
%     width=1.1\textwidth,
%     height=0.94\textwidth,
%     every axis plot/.append style={thick},
%     tick label style={font=\small},
%     label style={font=\small},
%     legend style={font=\small,
%        at={(0.55,0.29)}, % (x,y) 0~1 范围
%         anchor=south west
%     },
% ]

% % DuckDB
% \addplot[color=blue, mark=o, mark size=1pt]
% coordinates {
%     (0,0)(50,6447)(100,12894)(150,19341)(200,25788)(250,32235)
%     (300,38680)(350,39170)(400,39660)(450,40150)(500,40640)(550,40996)
%     (600,39996)(650,40090)(700,40184)(750,40278)(800,40321)(850,40321)
%     (900,40321)(950,40380)(1000,40439)(1050,40498)(1100,40557)(1150,40576)
%     (1200,40576)(1250,40610)(1300,40644)(1350,40678)(1400,40712)(1450,40777)
%     (1500,40777)
% };
% \addlegendentry{DuckDB}

% % PostgreSQL
% \addplot[color=red, mark=square, mark size=1pt]
% coordinates {
%     (0,0)(50,38850)(100,44685)(150,46520)(200,46608)(250,46696)
%     (300,46819)(350,46881)(400,46942)(450,46957)(500,46972)(550,47016)
%     (600,47060)(650,47064)(700,47068)(750,47074)(800,47080)(850,47085)
%     (900,47090)(950,47094)(1000,47097)(1050,47108)(1100,47118)(1150,47119)
%     (1200,47120)(1250,47128)(1300,47135)(1350,47142)(1400,47149)(1450,47149)
%     (1500,47149)
% };
% \addlegendentry{PostgreSQL}

% % SQLite
% \addplot[color=green!60!black, mark=triangle, mark size=1.5pt]
% coordinates {
%     (0,0)(50,9166)(100,10583)(150,10999)(200,11011)(250,11023)
%     (300,11024)(350,11025)(400,11028)(450,11031)(500,11032)(550,11033)
%     (600,11039)(650,11042)(700,11044)(750,11045)(800,11045)(850,11047)
%     (900,11049)(950,11050)(1000,11051)(1050,11052)(1100,11053)(1150,11053)
%     (1200,11054)(1250,11056)(1300,11057)(1350,11062)(1400,11066)(1450,11068)
%     (1500,11069)
% };
% \addlegendentry{SQLite}

% \end{axis}
% \end{tikzpicture}
% \end{minipage}
% \caption{Line coverage plateau.}
% \label{fig:challenge_plateau}
% \end{subfigure}
\caption{Illustrating challenges in LLM-based DBMS test case generation with Qwen2.5-7B.}
\label{fig:challenges}
\end{figure}



\noindent \textbf{Our work.} To address these challenges, we propose \textbf{MIST}, an LL\textbf{M}-based test case generat\textbf{I}on framework for DBM\textbf{S} through Monte Carlo \textbf{T}ree search. MIST consists of two complementary stages: \textit{1) Feature-Guided Error-Driven Test Case Synthetization.} This stage constructs a feature tree from official documentation and hierarchically samples feature combinations to guide LLM generation. Execution errors are captured and fed back to subsequent prompts, improving syntactic validity and semantic diversity of synthesized test cases. \textit{2) Monte Carlo Tree Search-Based Test Case Mutation.} After coverage growth plateaus, this stage employs MCTS to iteratively select seed queries and apply mutation rules guided by coverage feedback from test execution. With elaborate mutation rules covering schema, data, and query operations, this stage systematically explores deeper execution paths to boost code coverage.

We evaluate MIST on three widely-used open-source DBMSs (DuckDB, PostgreSQL, and SQLite) with four LLMs of different parameter scales (Qwen2.5-7B~\cite{Qwen2.5}, Llama3.1-8B~\cite{llama3.1}, Qwen2.5-14B~\cite{Qwen2.5}, and Qwen2.5-32B~\cite{Qwen2.5}). Experimental results show that MIST improves the baseline approach by 23.9\% $\sim$ 69.17\% on DuckDB, 18.5\% $\sim$ 54.1\% on PostgreSQL, and 18.5\% $\sim$ 92.9\% on SQLite in terms of average code coverage. Moreover, MIST achieves the highest line coverage in the \textit{Optimizer} module, reaching 69.3\% in DuckDB and 63.4\% in PostgreSQL. These results demonstrate that MIST greatly enhances test effectiveness for industrial DBMS testing using only lightweight, locally deployed LLMs. The source code is released at \url{https://github.com/yujiachen99/DBMSTesting}.

\noindent\textbf{Contributions.} The main contributions of our paper are summarized below:

\begin{itemize}[leftmargin=*]
    \item We propose MIST, a novel framework that effectively generates high-quality test cases to improve code coverage for proprietary DBMSs using only lightweight, locally deployed LLMs in resource-constrained industrial environments.
    
    \item We introduce a hierarchical feature tree with error-driven feedback and an MCTS-based mutation engine that jointly optimizes test case generation to improve code coverage.
    
    \item Extensive experiments on three widely-used DBMSs with four LLMs demonstrate that MIST remarkably outperforms baseline models in code coverage and critical module coverage.
    \end{itemize}


% \noindent Database Management Systems (DBMSs) lie at the core of modern data analytics and data-driven applications, supporting a wide range of real-world domains such as education, finance, healthcare, and e-commerce~\cite{}. The reliability of DBMSs is particularly critical in real-world scenarios, for example, in banking, healthcare, or financial trading systems, where even a minor flaw in any component may cascade into severe failures, compromising data integrity, query correctness, or overall system availability. DBMS testing plays a vital role in ensuring the correctness and robustness of these systems. In particular, generating high-quality SQL test cases is a fundamental prerequisite for validating DBMS behavior before deployment and for preventing subtle correctness bugs from propagating into production.
% In light of this, numerous approaches have been proposed in the literature to improve DBMS testing \cite{}. Traditional fuzzing and random SQL generation techniques can be effective for a specific DBMS; however, there exist many DBMSs designed for different application scenarios, such as DuckDB for analytical workloads \cite{}, MySQL and PostgreSQL for general-purpose OLTP systems \cite{}, and SQLite for lightweight embedded environments \cite{}.
% Adapting these methods to different DBMS dialects requires substantial manual effort, such as crafting grammar rules, maintaining dialect-specific operators, and supporting vendor-specific features, which severely limits their scalability~\cite{}.
% Recent advances in large language models (LLMs) provide a promising new opportunity to automate SQL test generation by leveraging knowledge distilled from documentation and real examples.
% Figure \ref{fig:intro_case} illustrates a representative example of using an LLM to directly generate SQL test cases.
% However, applying LLMs in real-world DBMS testing  still faces two fundamental challenges:

% \begin{itemize}[topsep=1mm, partopsep=0pt, itemsep=0pt, leftmargin=10pt]

% \item  \textbf{Challenge 1: Limited model adaptability in proprietary and resource-constrained industrial DBMS environments.}

% Limited semantic correctness and richness in local LLMs.} 
% Due to the limited capacity of locally deployed lightweight LLMs, they often generate extremely simple SQL queries that lack meaningful testing intent (see Figure~\ref{fig:challenges}(a)), or produce queries that violate DBMS-specific dialect rules (see Figure~\ref{fig:challenges}(b)). 
% These issues arise because such models lack sufficient domain knowledge and have incomplete exposure to proprietary SQL variants, making them prone to producing both semantically simplistic SQL and syntactically incorrect test cases.
% In industrial environments, this challenge becomes even more severe: strict security and privacy requirements typically restrict engineers to local LLMs, often 70B parameters or smaller, that cannot rely on powerful cloud-based closed-source models for assistance. As a result, ensuring syntactic correctness and semantic richness remains difficult, even when LLMs are employed for SQL generation.




% \item \textbf{Challenge 2: Limited diversity in LLM-generated test cases.}
% Simply generating a large number of SQL queries with an LLM does not translate into high code coverage, because many of the generated queries are repetitive, semantically similar, and exercise only shallow or common execution paths.
% As shown in Figure~\ref{fig:challenges}(c), LLM-generated test cases often exhibit limited structural and functional diversity, leading to insufficient coverage in practice. Consequently, they struggle to explore deeper or less common execution behaviors, causing the overall test coverage to quickly plateau.
% \end{itemize}



% To address these challenges, we propose MIST, a new framework that enables efficient, dialect-aware, and coverage-driven SQL test generation using only lightweight, locally deployed LLMs. MIST consists of two complementary stages. 
% {\it Stage I: Feature-Guided Error-Driven Test Case Synthetization} constructs a hierarchical multi-DBMS feature tree from official documentation and performs hierarchical path sampling with diversity enhancement to guide the LLM in producing syntactically valid and semantically rich SQL queries. Execution errors are captured and stored in a memory module, which then provides iterative feedback to help the LLM refine subsequent generations. 
% {\it Stage II: Monte Carlo Tree Search-Based Test Case Mutation} tackles the challenge of coverage stagnation by formulating SQL mutation as a tree-structured search problem. MIST employs MCTS to jointly optimize seed-selection and rule-application decisions, supported by a curated collection of 135 dialect-aware mutation rules. An early-termination strategy further improves search efficiency by pruning saturated seeds and prioritizing promising mutation trajectories.
% By unifying feature-guided LLM generation with coverage-guided search, MIST achieves both high dialect correctness and deep execution-path exploration, leading to significant improvements in code coverage for industrial DBMS testing. In doing so, MIST directly addresses the two key challenges identified above.

% In summary, our main contributions are as follows.
% \begin{itemize}
% \item
% We propose MIST, a two-stage dialect-aware and coverage-oriented SQL test generation system that effectively handles DBMS-specific SQL dialects while enabling deep execution-path exploration.
% \item
% We formulate SQL mutation as a search problem and introduce the first MCTS-based mutation engine that jointly optimizes seed selection and mutation-rule application, significantly improving the code coverage of the generated test cases.

% \item 
% Experiments on real industrial DBMSs show that MIST significantly outperforms existing fuzzing and LLM-based baselines in both code coverage and error-discovery capability. The source code is released at \url{https://anonymous.4open.science/r/DBMSTesting-DB84/}
% \end{itemize}

{\bf Outline.} 
We present the necessary background in \Cref{sec:background}.
\Cref{sec:overview} presents the overall framework of MIST, followed by detailed introduction. The experimental setup and results are provided in \Cref{sec:exp_setup} and \Cref{sec:result}, respectively.
\Cref{sec:discuss} offers a case study illustrating why MIST works and discusses threats to validity. Finally, we review related work in \Cref{sec:related} and conclude the paper in \Cref{sec:conclusions}.